% !PRINTED 139
\weberpage{139}
%<<CONT
1 with 2 and 4 with 5 are interchanged, into the first, and in the same way the three other assumptions about $b$ can be reduced to one another. There remain, then, only two possibilities to investigate:
\[
\begin{gathered}
a = (1, 2, 3) (4, 5, 6), \quad b = (1, 2, 3, 4, 5) \\
a = (1, 2, 3) (4, 5, 6), \quad b = (1, 2, 3, 5, 4).
\end{gathered}
\]
Of these, however, the first is to be rejected, because in
\[
a^2 b = (1, 4, 6)
\]
it would yield a three-term cycle, which cannot occur. There remains, then, only the single possibility:
\[
a = (1, 2, 3) (4, 5, 6), \quad b = (1, 2, 3, 5, 4),
\]
and there is therefore, if one regards equivalent groups as not different, only \gesperrt{one} simple group of the 60th degree.

These two elements $a, b$ may be regarded as generating elements of the whole group, and by their repeated composition one can derive the whole group of the 60th degree in a great variety of ways. Besides the unit element we obtain the permutations of the 5th degree:
\[
\begin{gathered}
(2, 3, 6, 5, 4), \quad (1, 3, 5, 6, 4), \quad (1, 2, 5, 4, 6), \\
(1, 2, 6, 3, 5), \quad (1, 2, 4, 6, 3), \quad (1, 2, 3, 5, 4),
\end{gathered}
\]
which with their powers yield 24; further the permutations of the 3rd degree:
\[
\begin{array}{ll}
(1, 2, 3) (4, 5, 6), & (1, 3, 5) (2, 4, 6), \\
(1, 2, 4) (3, 6, 5), & (1, 3, 6) (2, 4, 5), \\
(1, 2, 5) (3, 6, 4), & (1, 4, 5) (2, 3, 6), \\
(1, 2, 6) (3, 4, 5), & (1, 4, 6) (2, 3, 5), \\
(1, 3, 4) (2, 6, 5), & (1, 5, 6) (2, 3, 4),
\end{array}
\]
which together with their squares yield 20. Finally one obtains 15 more permutations of the 2nd degree:
\[
\begin{array}{lll}
(1, 2) (3, 4), & (1, 2) (5, 6), & (3, 4) (5, 6), \\
(1, 3) (4, 5), & (1, 3) (2, 6), & (2, 6) (4, 5), \\
(1, 4) (2, 5), & (1, 4) (3, 6), & (2, 5) (3, 6), \\
(1, 5) (2, 3), & (1, 5) (4, 6), & (2, 3) (4, 6), \\
(1, 6) (2, 4), & (1, 6) (3, 5), & (2, 4) (3, 5),
\end{array}
\]
of which each three standing in one row form, together with the unit element, a group of the 4th degree.

That this group can also be represented by the permutations of five figures now follows already from the fact
%CONT>>
